\section{Provas de consistência}\label{sec:provasDeConsistencia}
Uma prova de consistência é uma prova metateórica de que se certa teoria não apresenta contradições, então outra teoria também não apresenta contradições.

\begin{metadefinition}
    Seja $T$ uma teoria de primeira ordem em uma linguagem $\mathcal L$.

    Dizemos que $T$ é inconsistente se existe uma sentença $\phi$ tal que $T\vdash \phi$ e $T\vdash \neg \phi$.

    Se $T$ não é inconsistente, dizemos que $T$ é consistente, e escrevemos $\Con(T)$.
\end{metadefinition}

Conforme vimos, o Princípio da Explosão implica que uma teoria inconsistente prova qualquer sentença (ver Metaproposição~\ref{mprop:principioDaExplosao}).
Tal fato remove quase que completamente o interesse matemático por teorias de primeira ordem inconsistentes (no sistema da lógica clássica).

Assim, seria útil termos um método para provar que certas teorias de primeira ordem são consistentes, ao menos com relação a outras teorias de primeira ordem.

Iniciaremos com uma proposição simples.

\begin{metaproposition}\label{mprop:consistenciaDeTeoriasMaisFracas}
    Sejam $T$ e $T'$ teorias de primeira ordem com léxicos $\mathcal L$ e $\mathcal L'$, respectivamente, tais que todo símbolo de $\mathcal L$ é um símbolo de $\mathcal L'$.
    Suponha que todo axioma de $T$ é um teorema de $T'$.
    
    Se $\Con(T')$, então $\Con(T)$.
\end{metaproposition}

\begin{proof}
    Procedemos pela contrapositiva.

    Suponha que $\Con(T)$ é falso.
    Então $T$ é inconsistente, e existe uma sentença $\phi$ tal que $T\vdash \phi$ e $T\vdash \neg \phi$.
    Pela Metaproposição~\ref{mprop:extensaoDeTeoria}, $T'$ prova todos os teoremas de $T$, e, portanto, $T'\vdash \phi$ e $T'\vdash \neg \phi$.
    Assim, $T'$ é inconsistente.
\end{proof}

Teorias conservativas preservam consistência.
Em particular, teorias obtidas por extensão por definição de uma teoria $T$ são conservativas sobre $T$.

\begin{metaproposition}
    Sejam $T$ e $T'$ teorias de primeira ordem tais que $T'$ é uma extensão conservativa de $T$.

    Então $T$ é inconsistente se, e somente se, $T'$ é inconsistente.

    Dessa forma, $\Con(T)$ se, e somente se, $\Con(T')$.
\end{metaproposition}
\begin{proof}
    Suponha que $T$ é inconsistente.
    Pela Metaproposição~\ref{mprop:consistenciaDeTeoriasMaisFracas}, $T'$ é inconsistente.

    Agora suponha que $T'$ é inconsistente.
    Temos, pelo Princípio da Explosão, que $T'\vdash \neg \forall x\, (x=x)$.
    Pela conservatividade, $T\vdash \neg \forall x\, (x=x)$.
    Por outro lado, $T\vdash \forall x\, (x=x)$, pois é um teorema de qualquer teoria de primeira ordem.
    Assim, $T$ é inconsistente.
\end{proof}

Uma importante forma de produzir provas de consistência em Teoria dos Conjuntos é por meio de examinar o universo dos conjuntos restrito a alguma classe definível.

\begin{metadefinition}
    Seja $T$ uma teoria de primeira ordem e $\mathcal L'$ um léxico para uma linguagem de primeira ordem.
    Uma \emph{interpretação} de $\mathcal L'$ em $T$ é uma extensão por definição $T^A$ de $T$ tal que a linguagem de $T^A$ é obtida a partir da linguagem de $T$ por meio da adição de:
    \begin{itemize}
        \item Um novo símbolo de relação unário $A$, introduzido na linguagem de $T$ via definição, de modo que $T^A\vdash \exists x\, A(x)$.
        \item Para cada símbolo de predicado $R$ de $\mathcal L'$ de aridade $n$, um símbolo de predicado $R^A$ de aridade $n$ introduzido na linguagem de $T$ via definição.
        \item Para cada símbolo de constante $c$ de $\mathcal L'$, um símbolo de constante $c^A$ introduzido na linguagem de $T$ via definição de modo que $T^A\vdash A(c^A)$.
        \item Para cada símbolo de função $f$ de $\mathcal L'$ de aridade $n\geq1$, um símbolo de função $f^A$ de aridade $n$ introduzido na linguagem de $T$ via definição de modo que
        \[
            T^A\vdash \forall x_1 \ldots \forall x_n\, (A(x_1)\wedge \cdots \wedge A(x_n)\rightarrow A(f^A(x_1, \dots, x_n))).
        \]
    \end{itemize}

    Para cada $\mathcal L'$-termo $t$, criamos o termo $t^A$ da linguagem de $T^A$ por recursão na metateoria, de modo que:
    \begin{itemize}
        \item Se $t$ é uma variável, digamos, $x$, então $t^A$ é a própria variável $x$.
        \item Se $t$ é uma constante $c$, então $t^A$ é a constante $c^A$.
        \item Se $t$ é da forma $f(t_1, \dots, t_n)$, então $t^A$ é o termo $f^A(t_1^A, \dots, t_n^A)$.
    \end{itemize}

    Para cada $\mathcal L'$-fórmula $\phi$, criamos a fórmula $\phi^A$ da linguagem de $T^A$ por recursão na metateoria, de modo que:
    \begin{itemize}
        \item Se $\phi$ é um símbolo de predicado $0$-ário $R$, então $\phi^A$ é a fórmula $R^A$.
        \item Se $\phi$ é da forma $R(t_1, \dots, t_n)$ em que $R$ é um símbolo de predicado de $\mathcal L'$ de aridade $n\geq 1$, então $\phi^A$ é a fórmula $R^A(t_1^A, \dots, t_n^A)$.
        \item Se $\phi$ é da forma $t_1=t_2$, então $\phi^A$ é a fórmula $t_1^A=t_2^A$.
        \item Se $\phi$ é da forma $\neg \psi$, então $\phi^A$ é a fórmula $\neg \psi^A$.
        \item Se $\phi$ é da forma $\psi\wedge \theta$, então $\phi^A$ é a fórmula $\psi^A\wedge \theta^A$.
        \item Se $\phi$ é da forma $\psi\vee \theta$, então $\phi^A$ é a fórmula $\psi^A\vee \theta^A$.
        \item Se $\phi$ é da forma $\psi\to \theta$, então $\phi^A$ é a fórmula $\psi^A\to \theta^A$.
        \item Se $\phi$ é da forma $\psi\leftrightarrow \theta$, então $\phi^A$ é a fórmula $\psi^A\leftrightarrow \theta^A$.
        \item Se $\phi$ é da forma $\forall x\,\psi$, então $\phi^A$ é a fórmula $\forall x (A(x)\to \psi^A)$.
        \item Se $\phi$ é da forma $\exists x\,\psi$, então $\phi^A$ é a fórmula $\exists x (A(x)\wedge \psi^A)$.
    \end{itemize}

    Se $S$ é uma teoria de primeira ordem de léxico $\mathcal L'$, dizemos que $T^A$ é uma \emph{interpretação} de $S$ em $T$ se $T^A$ é uma interpretação de $\mathcal L'$ em $T$ e se, para todo axioma $\sigma$ de $S$, $T^A\vdash \sigma^A$.

    Algumas convenções notacionais: muitas vezes, para abreviar fórmulas de modo mais intuitivo, adotamos as seguintes convenções.

    \begin{itemize}
        \item Escrevemos $x\in A$ em vez de $A(x)$.
        \item Escrevemos $\exists x\in A\,\psi$ em vez de $\exists x (A(x)\wedge \psi)$.
        \item Escrevemos $\forall x\in A\,\psi$ em vez de $\forall x (A(x)\to \psi)$.
    \end{itemize}
\end{metadefinition}

Como $T^A$ é uma extensão por definição de $T$, $T^A$ é uma extensão conservativa de $T$, e, portanto, $\Con(T)$ se, e somente se, $\Con(T^A)$.

Notemos que $\neg\forall x\,\psi$ é logicamente equivalente a $\exists x\,\neg\psi$, e que $\neg\exists x\,\psi$ é logicamente equivalente a $\forall x\,\neg\psi$.

Agora faremos algumas observações simples sobre interpretações de termos e fórmulas.

\begin{metalemma}
    Seja $T$ uma teoria de primeira ordem e seja $T^A$ uma interpretação de um léxico $\mathcal L'$ em $T$.

    Se $t$ é um $\mathcal L'$-termo, então as variáveis que ocorrem em $t$ são as mesmas que as variáveis que ocorrem em $t^A$.
    
    Se $\phi$ é uma $\mathcal L'$-fórmula, então as variáveis livres de $\phi$ são as mesmas que as variáveis livres de $\phi^A$.
\end{metalemma}
\begin{proof}
    Para a primeira afirmação, a demonstração é uma indução simples na complexidade de $t$.

    Se $t$ é uma variável, então $t^A$ é a própria variável $t$.

    Se $t$ é uma constante $c$, então $t^A$ é a constante $c^A$, e, portanto, nenhum dos dois possui variáveis.

    Se $t$ é da forma $f(t_1,\dots,t_n)$, então $t^A$ é da forma $f^A(t_1^A,\dots,t_n^A)$.
    Assim, as variáveis que ocorrem em $t$ são exatamente as variáveis que ocorrem em algum $t_i$, e as variáveis que ocorrem em $t^A$ são exatamente as variáveis que ocorrem em algum $t_i^A$.
    Pela hipótese indutiva, para cada $1\leq i\leq n$, as variáveis que ocorrem em $t_i$ são as mesmas que as variáveis que ocorrem em $t_i^A$.
    Portanto, as variáveis que ocorrem em $t$ são as mesmas que as variáveis que ocorrem em $t^A$.

    Já para a segunda afirmação, a demonstração é por indução na complexidade de $\phi$.

    Se $\phi$ é um símbolo de predicado $0$-ário, então nem $\phi$ nem $\phi^A$ possuem variáveis livres.

    Se $\phi$ é atômica da forma $R(s_1,\dots,s_n)$, com $n\geq 1$, então $\phi^A$ é $R^A(s_1^A,\dots,s_n^A)$.
    As variáveis livres de $\phi$ são exatamente as variáveis que ocorrem em algum $s_i$, e as variáveis livres de $\phi^A$ são exatamente as variáveis que ocorrem em algum $s_i^A$.
    Pela primeira afirmação, essas variáveis são as mesmas.

    Se $\phi$ é atômica da forma $s=t$, então $\phi^A$ é $s^A=t^A$.
    Pela primeira afirmação, as variáveis que ocorrem em $s$ e em $t$ são, respectivamente, as mesmas que ocorrem em $s^A$ e em $t^A$.
    Portanto, as variáveis livres de $\phi$ são as mesmas que as variáveis livres de $\phi^A$.

    Se $\phi$ é da forma $\neg\psi$, então $\phi^A$ é $\neg\psi^A$.
    Pela hipótese indutiva, as variáveis livres de $\psi$ são as mesmas que as variáveis livres de $\psi^A$.
    Portanto, as variáveis livres de $\phi$ são as mesmas que as variáveis livres de $\phi^A$.

    Se $\phi$ é da forma $\psi\square\theta$, com $\square$ um conectivo lógico binário, então $\phi^A$ é da forma $\psi^A\square\theta^A$.
    Pela hipótese indutiva, as variáveis livres de $\psi$ e de $\theta$ são, respectivamente, as mesmas que as variáveis livres de $\psi^A$ e de $\theta^A$.
    Portanto, as variáveis livres de $\phi$ são as mesmas que as variáveis livres de $\phi^A$.

    Se $\phi$ é da forma $\forall x\,\psi$, então $\phi^A$ é $\forall x\,(A(x)\to \psi^A)$.
    As variáveis livres de $\phi$ são as variáveis livres de $\psi$, exceto $x$.
    As variáveis livres de $\phi^A$ são as variáveis livres de $A(x)\to\psi^A$, exceto $x$; portanto, são as variáveis livres de $\psi^A$, exceto $x$.
    Pela hipótese indutiva, as variáveis livres de $\psi$ são as mesmas que as variáveis livres de $\psi^A$.
    Logo, as variáveis livres de $\phi$ são as mesmas que as variáveis livres de $\phi^A$.

    O caso em que $\phi$ é da forma $\exists x\,\psi$ é análogo.
\end{proof}

\begin{metalemma}
    Seja $T$ uma teoria de primeira ordem e seja $T^A$ uma interpretação de um léxico $\mathcal L'$ em $T$.
    Se $s$, $t$ são $\mathcal L'$-termos e $x$ é uma variável, então $(s[x/t])^A$ é $s^A[x/t^A]$.
\end{metalemma}
\begin{proof}
    A demonstração é por indução na complexidade de $s$.

    Se $s$ é uma variável $y$ diferente de $x$, então $s[x/t]$ é $y$.
    Logo, $(s[x/t])^A$ é $y$.
    Por outro lado, $s^A$ é $y$, e, portanto, $s^A[x/t^A]$ é $y$.

    Se $s$ é a variável $x$, então $s[x/t]$ é $t$.
    Logo, $(s[x/t])^A$ é $t^A$.
    Por outro lado, $s^A$ é $x$, e, portanto, $s^A[x/t^A]$ é $x[x/t^A]$, que é $t^A$.

    Se $s$ é uma constante $c$, então $s[x/t]$ é $c$.
    Logo, $(s[x/t])^A$ é $c^A$.
    Por outro lado, $s^A$ é $c^A$, e, como $c^A$ é uma constante, $s^A[x/t^A]$ é $c^A$.

    Finalmente, suponha que $s$ é da forma $f(s_1,\dots,s_n)$ e que a tese vale para $s_1,\dots,s_n$.
    Então
    \[
        (s[x/t])^A
        =
        f^A\bigl((s_1[x/t])^A,\dots,(s_n[x/t])^A\bigr).
    \]
    Pela hipótese indutiva, para cada $1\leq i\leq n$, temos $(s_i[x/t])^A=s_i^A[x/t^A]$.
    Portanto,
    \[
        (s[x/t])^A
        =
        f^A\bigl(s_1^A[x/t^A],\dots,s_n^A[x/t^A]\bigr).
    \]
    Pela definição de substituição em termos, o lado direito é exatamente
    \[
        f^A(s_1^A,\dots,s_n^A)[x/t^A],
    \]
    isto é, $s^A[x/t^A]$.
\end{proof}


\begin{metalemma}
    Seja $T$ uma teoria de primeira ordem e seja $T^A$ uma interpretação de um léxico $\mathcal L'$ em $T$.

    Se $\phi$ é uma $\mathcal L'$-fórmula, então as variáveis livres de $\phi$ são as mesmas que as variáveis livres de $\phi^A$.

    Além disso, se $t$ é um $\mathcal L'$-termo, $x$ é uma variável, e se a substituição $\phi[x/t]$ é válida, então a substituição $\phi^A[x/t^A]$ é válida e $(\phi[x/t])^A$ é $\phi^A[x/t^A]$.
\end{metalemma}
\begin{proof}
    Procedemos por indução na complexidade de $\phi$, provando simultaneamente as duas afirmações do enunciado.

    Se $\phi$ é um símbolo de predicado $0$-ário, então nem $\phi$ nem $\phi^A$ possuem variáveis livres.
    Além disso, $\phi[x/t]$ é $\phi$ e $\phi^A[x/t^A]$ é $\phi^A$.
    Logo, as duas afirmações são imediatas.

    Se $\phi$ é atômica da forma $R(s_1,\dots,s_n)$, com $n\geq 1$, então $\phi^A$ é $R^A(s_1^A,\dots,s_n^A)$.
    Pela preservação de variáveis em termos, as variáveis que ocorrem em cada $s_i$ são as mesmas que ocorrem em $s_i^A$.
    Portanto, as variáveis livres de $\phi$ são as mesmas que as variáveis livres de $\phi^A$.

    Como fórmulas atômicas não possuem quantificadores, a substituição de $x$ por $t^A$ em $\phi^A$ é válida.
    Pela compatibilidade da interpretação com substituições em termos, para cada $1\leq i\leq n$, $(s_i[x/t])^A$ é $s_i^A[x/t^A]$.
    Assim, $(\phi[x/t])^A$ é $R^A(s_1^A[x/t^A],\dots,s_n^A[x/t^A])$, que é exatamente $\phi^A[x/t^A]$.
  
    Se $\phi$ é atômica da forma $s_1=s_2$, o argumento é análogo.

    Se $\phi$ é da forma $\neg\theta$, então $\phi^A$ é $\neg\theta^A$.
    Pela hipótese indutiva, as variáveis livres de $\theta$ são as mesmas que as variáveis livres de $\theta^A$.
    Como a negação não altera variáveis livres, as variáveis livres de $\phi$ são as mesmas que as variáveis livres de $\phi^A$.

    Suponha agora que $\phi[x/t]$ é válida.
    Então $\theta[x/t]$ é válida.
    Pela hipótese indutiva, a substituição de $x$ por $t^A$ em $\theta^A$ é válida e $(\theta[x/t])^A$ é $\theta^A[x/t^A]$.
    Logo, a substituição de $x$ por $t^A$ em $\neg\theta^A$ é válida, e $((\neg\theta)[x/t])^A$ é $\neg(\theta^A[x/t^A])$,
    que é exatamente $(\neg\theta^A)[x/t^A]$.

    Se $\phi$ é da forma $\theta\square\chi$, com $\square$ um conectivo lógico binário, então $\phi^A$ é $\theta^A\square\chi^A$.
    Pela hipótese indutiva, as variáveis livres de $\theta$ e de $\chi$ são, respectivamente, as mesmas que as variáveis livres de $\theta^A$ e de $\chi^A$.
    Portanto, as variáveis livres de $\phi$ são as mesmas que as variáveis livres de $\phi^A$.

    Suponha agora que $\phi[x/t]$ é válida.
    Então $\theta[x/t]$ e $\chi[x/t]$ são válidas.
    Pela hipótese indutiva, a substituição de $x$ por $t^A$ é válida em $\theta^A$ e em $\chi^A$, e temos
    \[
        (\theta[x/t])^A \text{ é } \theta^A[x/t^A]
        \qquad\text{e}\qquad
        (\chi[x/t])^A \text{ é } \chi^A[x/t^A].
    \]
    Logo, a substituição de $x$ por $t^A$ em $\theta^A\square\chi^A$ é válida, e $((\theta\square\chi)[x/t])^A$
    é $\theta^A[x/t^A]\square\chi^A[x/t^A]$, que é exatamente $(\theta^A\square\chi^A)[x/t^A]$.

    Resta tratar os quantificadores.
    Seja $\phi$ da forma $Qy\,\theta$, em que $Q$ é $\forall$ ou $\exists$.
    Escrevamos $B_y(\eta)$ para $A(y)\to\eta$, se $Q$ é $\forall$, e para $A(y)\wedge\eta$, se $Q$ é $\exists$.
    Assim, $\phi^A$ é da forma
    \[
        Qy\,B_y(\theta^A).
    \]

    Pela hipótese indutiva, as variáveis livres de $\theta$ são as mesmas que as variáveis livres de $\theta^A$.
    As variáveis livres de $\phi$ são as variáveis livres de $\theta$, exceto $y$.
    Já as variáveis livres de $\phi^A$ são as variáveis livres de $B_y(\theta^A)$, exceto $y$.
    Como as variáveis livres de $B_y(\theta^A)$ são $y$ juntamente com as variáveis livres de $\theta^A$, segue que as variáveis livres de $\phi^A$ são as variáveis livres de $\theta^A$, exceto $y$.
    Portanto, as variáveis livres de $\phi$ são as mesmas que as variáveis livres de $\phi^A$.

    Suponha agora que $\phi[x/t]$ é válida.

    Primeiro, suponha que $y$ é diferente de $x$.
    Então $x$ é substituível por $t$ em $\theta$ e, além disso, ou $x$ não ocorre livre em $\theta$, ou $y$ não ocorre em $t$.
    Pela hipótese indutiva, $x$ é substituível por $t^A$ em $\theta^A$, e $(\theta[x/t])^A$ é $\theta^A[x/t^A]$.

    Pela preservação de variáveis em termos, se $y$ não ocorre em $t$, então $y$ não ocorre em $t^A$.
    Pela preservação de variáveis livres já demonstrada para $\theta$, se $x$ não ocorre livre em $\theta$, então $x$ não ocorre livre em $\theta^A$.
    Como $y$ é diferente de $x$, a variável $x$ não ocorre livre em $A(y)$.
    Logo, $x$ é substituível por $t^A$ em $B_y(\theta^A)$ e a condição referente ao quantificador inicial $Qy$ também é satisfeita.
    Portanto, a substituição de $x$ por $t^A$ em $\phi^A$ é válida.

    Além disso, $((Qy\,\theta)[x/t])^A$ é $(Qy\,\theta[x/t])^A$, que é $Qy\,B_y((\theta[x/t])^A)$.
    Pela hipótese indutiva, essa fórmula é $Qy\,B_y(\theta^A[x/t^A])$.
    Como $y$ é diferente de $x$, esta última fórmula é exatamente $(Qy\,B_y(\theta^A))[x/t^A]$,
    isto é, $\phi^A[x/t^A]$.

    Finalmente, suponha que $y$ é $x$.
    Então $\phi[x/t]$ é a própria $\phi$.
    Por outro lado, $\phi^A$ é
    \[
        Qx\,B_x(\theta^A).
    \]
    A variável $x$ não ocorre livre em $\phi^A$, pois todas as ocorrências livres de $x$ em $B_x(\theta^A)$ ficam no escopo do quantificador inicial $Qx$.
    Portanto, a substituição de $x$ por $t^A$ em $\phi^A$ é válida e $\phi^A[x/t^A]$ é a própria $\phi^A$.
    Assim, $(\phi[x/t])^A$ é $\phi^A[x/t^A]$.
\end{proof}

\begin{metalemma}\label{mlem:termosInterpretadosPertencemAA}
    Seja $T^A$ uma interpretação de um léxico $\mathcal L'$ em $T$.
    Seja $t$ um $\mathcal L'$-termo, e sejam $x_1,\dots,x_n$ variáveis que contêm todas as variáveis que ocorrem em $t$.
    Então
    \[
        T^A\vdash \forall x_1\in A\cdots\forall x_n\in A\,A(t^A).
    \]
\end{metalemma}
\begin{proof}
    A demonstração é por indução na complexidade de $t$.

    Se $t$ é uma variável, digamos, $x$, basta mostrar que $T^A\Vdash \forall x\in A\,(A(x)\rightarrow A(x))$, o que é trivial.

    Se $t$ é uma constante $c$, basta notar que $T^A\vdash A(c^A)$ pela definição de interpretação.

    Suponha agora que $t$ é da forma $f(t_1,\dots,t_m)$ e que as variáveis de $t_1,\dots,t_m$ estão entre $x_1,\dots,x_n$, variáveis distintas, e que vale a hipótese indutiva de que para cada $1\leq i\leq m$,
    \[
        T^A\vdash \forall x_1\in A\cdots\forall x_n\in A\,A(t_i^A).
    \]

    Escreva $t_i(x_1,\dots,x_n)$ para indicar que as variáveis de $t_i$ estão entre $x_1,\dots,x_n$ e utilizar substituições.

    Devemos ver que
    \[
        T^A\vdash \forall x_1\in A\cdots\forall x_n\in A\,A(f^A(t_1^A,\dots,t_m^A)).
    \]
    Trabalhando em $T^A$, fixe $a_1, \dots, a_n$ constantes para $x_1, \dots, x_n$.
    Basta ver que
    \[A(f^A(t_1^A(a_1,\dots,a_n),\dots,t_m^A(a_1,\dots,a_n))).
    \]
    Pela definição de interpretação, $\forall y_1\cdots\forall y_m\,
        \bigl(A(y_1)\wedge\cdots\wedge A(y_m)\to
        A(f^A(y_1,\dots,y_m))\bigr)$.
    Instanciando essa fórmula com os termos $t_1^A(a_1,\dots,a_n),\dots,t_m^A(a_1,\dots,a_n)$, obtemos
    \[
        A(t_1^A(a_1,\dots,a_n))\wedge\cdots\wedge A(t_m^A(a_1,\dots,a_n))\to
        A(f^A(t_1^A(a_1,\dots,a_n),\dots,t_m^A(a_1,\dots,a_n))).
    \]
    Por hipótese, para cada $1\leq i\leq m$, $A(t_i^A(a_1,\dots,a_n))$.
    Logo, $A(f^A(t_1^A(a_1,\dots,a_n),\dots,t_m^A(a_1,\dots,a_n)))$.
\end{proof}

\begin{metalemma}[Axiomas lógicos relativizados]\label{mlem:axiomasLogicosInterpretacao}
    Sejam $S$ e $T$ teorias de primeira ordem e seja $T^A$ uma interpretação de $S$ em $T$.
    Seja $\alpha$ um axioma lógico na linguagem de $S$.
    Se $x_1,\dots,x_n$ são exatamente as variáveis livres de $\alpha$, então
    \[
        T^A\vdash \forall x_1\in A\cdots \forall x_n\in A\,\alpha^A(x_1, \dots, x_n).
    \]
\end{metalemma}

\begin{proof}
    Seja $\alpha$ um axioma lógico da linguagem de $S$.

    Primeiro, considere os axiomas FO1--FO13.
    Se $\alpha$ é uma instância de um desses esquemas, então $\alpha^A$ é uma instância do mesmo esquema.

    O caso FO16 também é imediato.
    Se $\alpha$ é $x_1=x_1$, então $\alpha^A$ é $a_1=a_1$, que é um teorema.

    Considere agora FO17.
    Suponha que $\alpha$ é
    \[
        \bigl(x_1=y_1\wedge\cdots\wedge x_m=y_m\bigr)
        \to
        f(x_1,\dots,x_m)=f(y_1,\dots,y_m),
    \]

 
    em que $f$ é um símbolo funcional de aridade $m>0$ da linguagem de $S$.

     Nesse caso, consideremos que os nomes das constantes fixadas para $x_1, \dots, x_n$ e para $y_1, \dots, y_m$ são $
     Então $\alpha^A$ é
    \[
        \bigl(x_1=y_1\wedge\cdots\wedge x_m=y_m\bigr)
        \to
        f^A(x_1,\dots,x_m)=f^A(y_1,\dots,y_m),
    \]
    que é uma instância de FO17 na linguagem de $T^A$.
    Logo, $T^A\vdash \alpha^A$, e a conclusão relativizada segue por introdução dos quantificadores restritos.

    O caso FO18 é análogo.
    Se $R$ é um símbolo de predicado não lógico da linguagem de $S$, então a interpretação de uma instância de FO18 para $R$ é a instância correspondente de FO18 para $R^A$ na linguagem de $T^A$.
    Se o predicado em questão é o símbolo lógico de igualdade, a tradução não altera o símbolo de igualdade, e obtemos novamente uma instância de FO18 na linguagem de $T^A$.

    Resta tratar FO14 e FO15.

    Suponha primeiro que $\alpha$ é uma instância de FO14, isto é,
    \[
        \alpha \quad\text{é}\quad
        \forall z\,\varphi\to \varphi[z/t],
    \]
    em que $z$ é substituível por $t$ em $\varphi$.
    Pelo metalema sobre substituições e interpretação, a substituição de $z$ por $t^A$ em $\varphi^A$ é válida e
    \[
        (\varphi[z/t])^A
    \]
    é
    \[
        \varphi^A[z/t^A].
    \]
    Assim, $\alpha^A$ é
    \[
        \forall z(A(z)\to\varphi^A)\to \varphi^A[z/t^A].
    \]

    Se $z$ ocorre livre em $\varphi$, então todas as variáveis de $t$ ocorrem livres em $\alpha$.
    Pelo Metalema~\ref{mlem:termosInterpretadosPertencemAA}, sob as hipóteses $x_1\in A,\dots,x_n\in A$, obtemos $A(t^A)$.
    Por FO14 em $T^A$, aplicado à fórmula $A(z)\to\varphi^A$ e ao termo $t^A$, temos
    \[
        T^A\vdash
        \forall z(A(z)\to\varphi^A)
        \to
        \bigl(A(t^A)\to\varphi^A[z/t^A]\bigr).
    \]
    Usando proposicionalmente $A(t^A)$, obtemos
    \[
        \forall z(A(z)\to\varphi^A)
        \to
        \varphi^A[z/t^A].
    \]
    Portanto,
    \[
        T^A\vdash
        \forall x_1\in A\cdots\forall x_n\in A\,\alpha^A.
    \]

    Se $z$ não ocorre livre em $\varphi$, então $\varphi[z/t]$ é a própria $\varphi$.
    Nesse caso, $\alpha^A$ é
    \[
        \forall z(A(z)\to\varphi^A)\to \varphi^A.
    \]
    Como $T^A\vdash \exists z\,A(z)$ pela definição de interpretação, a conclusão segue de FO14, da regra de introdução do quantificador existencial e de modus ponens.
    Com efeito, escolhendo uma variável $w$ que não ocorre livre em $\varphi^A$, FO14 dá
    \[
        \forall z(A(z)\to\varphi^A)\to (A(w)\to\varphi^A).
    \]
    Por raciocínio proposicional,
    \[
        A(w)\to\bigl(\forall z(A(z)\to\varphi^A)\to\varphi^A\bigr).
    \]
    Como $w$ não ocorre livre em $\forall z(A(z)\to\varphi^A)\to\varphi^A$, pela regra de introdução do quantificador existencial,
    \[
        \exists w\,A(w)\to\bigl(\forall z(A(z)\to\varphi^A)\to\varphi^A\bigr).
    \]
    Como $T^A\vdash \exists w\,A(w)$, obtemos $\alpha^A$.
    Relativizando as variáveis livres $x_1,\dots,x_n$, obtemos a conclusão desejada.

    Finalmente, suponha que $\alpha$ é uma instância de FO15, isto é,
    \[
        \alpha \quad\text{é}\quad
        \varphi[z/t]\to \exists z\,\varphi,
    \]
    em que $z$ é substituível por $t$ em $\varphi$.
    Pelo metalema sobre substituições e interpretação, $\alpha^A$ é
    \[
        \varphi^A[z/t^A]\to \exists z(A(z)\wedge\varphi^A).
    \]

    Se $z$ ocorre livre em $\varphi$, então todas as variáveis de $t$ ocorrem livres em $\alpha$.
    Pelo Metalema~\ref{mlem:termosInterpretadosPertencemAA}, sob as hipóteses $x_1\in A,\dots,x_n\in A$, obtemos $A(t^A)$.
    Pela introdução de $\wedge$,
    \[
        A(t^A)\wedge \varphi^A[z/t^A].
    \]
    Como
    \[
        A(t^A)\wedge \varphi^A[z/t^A]
    \]
    é
    \[
        (A(z)\wedge\varphi^A)[z/t^A],
    \]
    FO15 em $T^A$, aplicado à fórmula $A(z)\wedge\varphi^A$ e ao termo $t^A$, dá
    \[
        A(t^A)\wedge \varphi^A[z/t^A]
        \to
        \exists z(A(z)\wedge\varphi^A).
    \]
    Portanto, proposicionalmente,
    \[
        \varphi^A[z/t^A]\to \exists z(A(z)\wedge\varphi^A).
    \]
    Relativizando as variáveis livres, obtemos
    \[
        T^A\vdash
        \forall x_1\in A\cdots\forall x_n\in A\,\alpha^A.
    \]

    Se $z$ não ocorre livre em $\varphi$, então $\varphi[z/t]$ é a própria $\varphi$, e $\alpha^A$ é
    \[
        \varphi^A\to \exists z(A(z)\wedge\varphi^A).
    \]
    Como $T^A\vdash \exists w\,A(w)$ para uma variável $w$ que não ocorre livre em $\varphi^A$, obtemos, por FO15 aplicado à fórmula $A(z)\wedge\varphi^A$ e ao termo $w$,
    \[
        A(w)\wedge\varphi^A \to \exists z(A(z)\wedge\varphi^A).
    \]
    Logo, proposicionalmente,
    \[
        A(w)\to\bigl(\varphi^A\to\exists z(A(z)\wedge\varphi^A)\bigr).
    \]
    Pela regra de introdução do quantificador existencial, segue que
    \[
        \exists w\,A(w)\to
        \bigl(\varphi^A\to\exists z(A(z)\wedge\varphi^A)\bigr).
    \]
    Como $T^A\vdash \exists w\,A(w)$, obtemos $\alpha^A$.
    Relativizando as variáveis livres, concluímos novamente
    \[
        T^A\vdash
        \forall x_1\in A\cdots\forall x_n\in A\,\alpha^A.
    \]

    Isso encerra todos os casos de axiomas lógicos FO1--FO18.
\end{proof}